\chapter[CONCLUSION]{Conclusion}
\label{ch:conclusion}

This proposal sets out a contactless framework that detects corrosion on metal
surfaces by fusing RGB, thermal and LiDAR data, and measures the corroded area
and metal loss of each detected region from the LiDAR.

The case for the work rests on the limits of a single camera. Fitness for
service depends on how deep the metal loss goes as well as how much area it
covers, yet a camera records appearance only. The review in
Chapter~\ref{ch:literature-review} found that, of 33 corrosion-sensing studies,
16 rely on RGB images alone, only one adds thermal imaging, and none fuses RGB,
thermal and LiDAR data; nor has any study measured what each of these sensors
contributes to corrosion detection under controlled conditions.

The three objectives address this gap: a pipeline that registers the three
sensors into a common reference frame; an early-fusion YOLO11 detector with the
corroded area and metal loss of each region measured from the LiDAR; and an
ablation, with a blank-channel control and a decision rule fixed in advance,
that measures what each sensor adds. A preliminary test on the public MID-1K
dataset showed that the pipeline and the protocol run end to end. Early fusion
did not clear the decision rule, and loose RGB--thermal alignment cost the
thermal images nearly a third of their score, which confirms that registration
must come first.

The expected contributions are a registered RGB--thermal--LiDAR corrosion
dataset, to be released if the owners of the scanned structures permit, and a
controlled measure of which sensors are worth deploying for corrosion detection.
The main risk is data acquisition; if field access is not secured in time, the
work will be evaluated on controlled corrosion specimens and the resulting limit
on generalization will be reported.
